\section*{Chapitre \ref{ch:m2:mortgages-and-agencies} --- Prêts hypothécaires et agences}

\begin{solution}{exo:m2:mortgages-and-agencies:1}
2~\% ($10 \times 0.2\%$)~; 6~\% ($1.5 \times 20 \times 0.2\%$)~; 12~\% ($2 \times 6\%$,
la rampe étant terminée au bout de trente mois).
\end{solution}

\begin{solution}{exo:m2:mortgages-and-agencies:2}
$1 - 0.94^{1/12} = 0.5143\%$ par mois~; 257\,151~USD se remboursent par anticipation.
\end{solution}

\begin{solution}{exo:m2:mortgages-and-agencies:3}
Agence, coupon, échéance, prix, valeur nominale, mois de règlement-livraison. Le vendeur
livre les pools les moins valorisés qui respectent les critères de bonne livraison, ceux
dont on attend le \omterm{def:m2:mortgages-and-agencies:cpr}{remboursement anticipé} de la façon la plus défavorable au détenteur. Les
pools qui valent davantage, parce que leurs emprunteurs devraient rembourser plus lentement, sont
mis de côté et vendus comme \omterm{def:m2:mortgages-and-agencies:tba}{specified pools}, moyennant un pay-up.
\end{solution}

\begin{solution}{exo:m2:mortgages-and-agencies:4}
5.90~\%. Sans baisse, le roll abandonne le coupon du mois (6~\% par an au pair)
et le remboursement, restitué au pair plutôt qu'à 100.50~; financer 100.50 de
trésorerie à un taux un peu inférieur au coupon, voilà ce que cela coûte.
\end{solution}

\begin{solution}{exo:m2:mortgages-and-agencies:5}
À 4~\%, l'incitation est de deux points et demi, et le pool se rembourse à un CPR de
48~\%~: son principal revient vite au pair et ne peut pas être réinvesti à 6~\%.
Le pool à vitesse figée continue de verser 6~\% sur un encours qui s'amortit à la
vitesse à la monnaie, ce qui vaut bien plus à un taux d'actualisation de 4~\%.
\end{solution}

\begin{solution}{exo:m2:mortgages-and-agencies:6}
À 5.5~\%~: 1.77 ans et $-478$~; à 6~\%~: 4.73 et $-560$~; à 6.5~\%~: 6.51 et
$-93$. La \omterm{def:m2:mortgages-and-agencies:convexity}{convexité négative} est maximale près de la monnaie, là où la courbe en S est
la plus pentue et où une faible variation des taux modifie le plus les vitesses~; hors de la monnaie, les vitesses
réagissent à peine et le pool se comporte davantage comme une obligation.
\end{solution}

\begin{solution}{exo:m2:mortgages-and-agencies:7}
19.62, 14.73, 11.49, 7.73 et 5.77 ans.
\end{solution}

\begin{solution}{exo:m2:mortgages-and-agencies:8}
La garantie supprime le risque de crédit, pas le risque de \omterm{def:m2:mortgages-and-agencies:cpr}{remboursement anticipé}. Le
détenteur est vendeur de l'option qu'ont les emprunteurs de rembourser au pair, de sorte que les titres d'agence doivent rapporter davantage que
les Treasuries de maturité comparable pour rémunérer cette option, leurs flux de trésorerie
incertains et leur \omterm{def:m2:mortgages-and-agencies:convexity}{convexité négative}. Pour Fannie Mae et Freddie Mac, la
garantie est celle des entreprises, et non la garantie pleine et entière des
États-Unis qui soutient Ginnie Mae.
\end{solution}

\begin{solution}{pb:m2:mortgages-and-agencies:1}
\textbf{1.} CPR 11.9~\%, prix 100.00, \omterm{def:m2:mortgages-and-agencies:convexity}{duration effective} 4.73 ans.
\textbf{2.} 4.73~millions de USD par point de base ($10^{10} \times 4.727 \times 10^{-4}$).
\textbf{3.} 73\,601~USD.
\textbf{4.} Environ 6.42~milliards de USD.
\textbf{5.} $-560$~: la duration se raccourcit quand les taux baissent et s'allonge
quand ils montent, de sorte que la couverture doit être négociée dans le sens de chaque mouvement.
\textbf{6.} CPR 21.3~\%, prix 101.61, \omterm{def:m2:mortgages-and-agencies:convexity}{duration effective} 1.77 ans.
\textbf{7.} 1.80~million de USD par point de base.
\textbf{8.} À hauteur de 2.93~millions de USD par point de base~: l'investisseur est
vendeur de trop de duration et doit la racheter.
\textbf{9.} Environ 3.97~milliards de USD, en recevant le fixe.
\textbf{10.} 6.51 ans~: la couverture devrait plutôt être augmentée.
\textbf{11.} Environ 1.4~\%.
\textbf{12.} Après une baisse, les hedgers reçoivent le fixe ou achètent des obligations, ce qui pousse les taux
encore plus bas~; après une hausse, ils paient le fixe ou vendent.
\textbf{13.} Les investisseurs ayant le besoin inverse, comme ceux qui paient le fixe
pour se couvrir contre la hausse des taux, et les dealers, qui font payer ce service~; quand ils sont rares,
les taux s'écartent davantage.
\textbf{14.} Les options achetées à l'avance apportent automatiquement la duration quand
les taux baissent, de sorte que l'investisseur négocie moins après le mouvement~; la convexité a été
transférée au vendeur de l'option.
\textbf{15.} Elle a acheté des pools et ne les a pas couverts, de sorte que moins de la
\omterm{def:m2:mortgages-and-agencies:convexity}{convexité négative} du marché était couverte de manière dynamique.
\textbf{16.} La courbe en S, sa pente et le spread entre le taux hypothécaire et le taux
d'actualisation sont des hypothèses~; les vitesses réelles dépendent aussi de l'âge des prêts, de leur taille, de la
solvabilité des emprunteurs et de la capacité des prêteurs. L'ordre de grandeur de la réponse est robuste,
ses chiffres, eux, ne le sont pas.
\textbf{17.} La rampe de type PSA~: les nouveaux emprunteurs refinancent rarement durant leurs premiers
mois, quelle que soit l'incitation.
\textbf{18.} Elle perd des revenus de commissions à mesure que l'encours se rembourse, de sorte que son actif perd de la valeur
lors d'une hausse des cours, et elle aussi doit acheter de la duration pour la couvrir.
\textbf{19.} Résultat à retenir~: \emph{la couverture de convexité} pour 10~milliards de USD après
une baisse de 50 points de base des taux consiste à recevoir le fixe sur des swaps à dix ans
pour environ 4.0~milliards de USD, la \omterm{def:m2:mortgages-and-agencies:convexity}{duration effective} passant de 4.73 à 1.77 ans.
\textbf{20.} Parce que les emprunteurs refinancent, leurs titres se raccourcissent, et les détenteurs couverts
doivent remplacer la duration qu'ils ont perdue.
\end{solution}

\begin{solution}{iq:m2:mortgages-and-agencies:1}
Les emprunteurs peuvent rembourser au pair à tout moment, de sorte que le détenteur est
vendeur d'une option d'achat sur le prêt. Quand les taux baissent, l'option entre dans la monnaie, les remboursements
s'accélèrent et le prix est plafonné près du pair~: la duration se raccourcit. Quand les taux
montent, les remboursements ralentissent et la duration s'allonge. Un prix qui monte moins
qu'il ne baisse est une \omterm{def:m2:mortgages-and-agencies:convexity}{convexité négative}.
\iqlookfor{l'option incorporée et les variations de duration qu'elle provoque.}
\end{solution}

\begin{solution}{iq:m2:mortgages-and-agencies:2}
Une transaction à terme sur des \omterm{def:m2:mortgages-and-agencies:mbs}{pass-through} d'agence, définie par six paramètres (agence, coupon,
échéance, prix, nominal, mois de règlement-livraison), les pools étant désignés deux jours avant
le règlement-livraison. Elle regroupe un million de titres hétérogènes en quelques
contrats fongibles, ce qui les rend liquides et peu coûteux à échanger~; les pools aux
caractéristiques recherchées s'échangent séparément sous forme de \omterm{def:m2:mortgages-and-agencies:tba}{specified pools}.
\iqlookfor{la fongibilité, et sa conséquence sur le \omterm{def:m2:bond-futures:ctd}{moins cher à livrer} (CTD).}
\end{solution}

\begin{solution}{iq:m2:mortgages-and-agencies:3}
Vendre le TBA de l'échéance la plus proche et acheter celle de l'échéance suivante~: le
vendeur du roll renonce au coupon et au remboursement pendant un mois et garde la trésorerie. La baisse de prix implique un
taux de financement implicite~; lorsqu'il est inférieur au repo, le roll est special. Si l'on possède les titres,
il faut les faire rouler et investir la trésorerie~; si l'on a besoin des pools de
l'échéance la plus proche, il faut s'attendre à les payer plus cher.
\iqlookfor{le taux de financement implicite et la comparaison avec le repo.}
\end{solution}

\begin{solution}{iq:m2:mortgages-and-agencies:4}
Séparer les facteurs~: ventes de logements (ancienneté, saison de l'année), refinancement
en fonction de l'incitation, de l'épuisement (burnout), de la taille du prêt, de la solvabilité et de la
capacité des prêteurs~; calibrer sur l'historique au niveau du prêt ou du pool~; valider hors échantillon,
en particulier à travers les vagues de refinancement passées, et sur la dispersion des
vitesses entre les pools, et non sur leur seule moyenne. Vérifier que les durations du modèle couvrent
bien en pratique.
\iqlookfor{les facteurs, la validation hors échantillon et un test de couverture.}
\end{solution}

\begin{solution}{iq:m2:mortgages-and-agencies:5}
Après une hausse des cours, les durations hypothécaires se raccourcissent et les hedgers reçoivent le fixe ou achètent
des Treasuries~; après une baisse des cours, ils paient le fixe ou vendent. Les flux vont dans le sens
du mouvement, ce qui ajoute de la dynamique et fait monter la volatilité réalisée et la volatilité implicite,
surtout lorsqu'une grande partie du marché est proche de la monnaie. Les avoirs des banques
centrales et la couverture par options atténuent l'effet.
\iqlookfor{le sens, l'ampleur et ce qui l'amortit.}
\end{solution}

\begin{solution}{iq:m2:mortgages-and-agencies:6}
Regrouper les pools selon les caractéristiques qui déterminent les vitesses, ou utiliser le modèle de \omterm{def:m2:mortgages-and-agencies:cpr}{remboursement anticipé}
vectorisé sur les pools~; réutiliser un jeu commun de scénarios de taux (hausse, baisse et scénario central,
ou des trajectoires simulées) et calculer les flux une seule fois par scénario~; paralléliser entre
les pools~; mettre en cache les résultats des pools dont les entrées n'ont pas changé~; ne recalculer que les
parties dépendant du marché en cours de journée. Rapprocher les résultats d'un calcul complet, plus lent, effectué chaque nuit,
et suivre les durations du modèle au regard de la performance réalisée des couvertures.
\iqlookfor{vectorisation, réutilisation des scénarios, mise en cache et rapprochement.}
\end{solution}
